For 2D image understanding, pixels are placed in regular grids and can be processed with classical convolution. In contrast, 3D point clouds are unordered and scattered in 3D space: they are essentially sets. Learning-based approaches to processing 3D point clouds can be classified into the following types: projection-based, voxel-based, and point-based networks.

\mypara{Projection-based networks.}
For processing irregular inputs like point clouds, an intuitive way is to transform irregular representations to regular ones. Considering the success of 2D CNNs, some approaches~\cite{su15mvcnn,li2016vehicle,chen2017multi,kanezaki2018rotationnet,lang2019pointpillars} adopt multi-view projection, where 3D point clouds are projected into various image planes. Then 2D CNNs are used to extract feature representations in these image planes, followed by multi-view feature fusion to form the final output representations. In a related approach, TangentConv~\cite{tatarchenko2018tangent} projects local surface geometry onto a tangent plane at every point, forming tangent images that can be processed by 2D convolution. However, this approach heavily relies on tangent estimation. In projection-based frameworks, the geometric information inside point clouds is collapsed during the projection stage. These approaches may also underutilize the sparsity of point clouds when forming dense pixel  grids on projection planes. The choice of projection planes may heavily influence recognition performance and occlusion in 3D may impede accuracy.

\mypara{Voxel-based networks.}
An alternative approach to transforming irregular point clouds to regular representations is 3D voxelization~\cite{maturana2015voxnet,song2016ssc}, followed by convolutions in 3D. When applied naively, this strategy can incur massive computation and memory costs due to the cubic growth in the number of voxels as a function of resolution. The solution is to take advantage of sparsity, as most voxels are usually unoccupied. For example, OctNet~\cite{riegler2017octnet} uses unbalanced octrees with hierarchical partitions. Approaches based on sparse convolutions, where the convolution kernel is only evaluated at occupied voxels, can further reduce computation and memory requirements~\cite{graham20183d,choy20194d}. These methods have demonstrated good accuracy but may still lose geometric detail due to quantization onto the voxel grid.

\mypara{Point-based networks.}
Rather than projecting or quantizing irregular point clouds onto regular grids in 2D or 3D, researchers have designed deep network structures that ingest point clouds directly, as sets embedded in continuous space.
PointNet~\cite{qi2017pointnet} utilizes permutation-invariant operators such as pointwise MLPs and pooling layers to aggregate features across a set. PointNet++~\cite{qi2017pointnet2} applies these ideas within a hierarchical spatial structure to increase sensitivity to local geometric layout.
Such models can benefit from efficient sampling of the point set, and a variety of sampling strategies have been developed~\cite{qi2017pointnet2,dovrat2019learning,wu2019pointconv,yang2019modeling,hu2020randla}.

A number of approaches connect the point set into a graph and conduct message passing on this graph. DGCNN~\cite{wang2019dgcnn} performs graph convolutions on kNN graphs. PointWeb~\cite{zhao2019pointweb} densely connects local neightborhoods.
ECC~\cite{simonovsky2017dynamic} uses dynamic edge-conditioned filters where convolution kernels are generated based on edges inside point clouds. SPG~\cite{landrieu2018spg} operates on a superpoint graph that represents contextual relationships. KCNet~\cite{shen2018mining} utilizes kernel correlation and graph pooling. Wang et al.~\cite{wang2018local} investigate the local spectral graph convolution. %operation.
GACNet~\cite{wang2019graph} employs graph attention convolution and HPEIN~\cite{jiang2019hpein} builds a hierarchical point-edge interaction architecture. DeepGCNs~\cite{li2019deepgcns} explore the advantages of depth in graph convolutional networks for 3D scene understanding.

A number of methods are based on continuous convolutions that apply directly to the 3D point set, with no quantization. PCCN~\cite{wang2018pccn} represents convolutional kernels as MLPs. SpiderCNN~\cite{xu2018spidercnn} defines kernel weights as a family of polynomial functions. Spherical CNN~\cite{esteves2018learning} designs spherical convolution to address the problem of 3D rotation equivariance. PointConv~\cite{wu2019pointconv} and KPConv~\cite{thomas2019kpconv} construct convolution weights based on the input coordinates. InterpCNN~\cite{mao2019interpolated} utilizes coordinates to interpolate pointwise kernel weights. PointCNN~\cite{li2018pointcnn} proposes to reorder the input unordered point clouds with special operators. Ummenhofer et al.~\cite{Ummenhofer2020} apply continuous convolutions to learn particle-based fluid dynamics.

\mypara{Transformer and self-attention.}
Transformer and self-attention models have revolutionized machine translation and natural language processing~\cite{vaswani2017attention,wu2019pay,Devlin2018,dai2019transformer,Yang2019xlnet}.
This has inspired the development of self-attention networks for 2D image recognition~\cite{hu2019local,ramachandran2019stand,zhao2020san,dosovitskiy2021image}. Hu et al.~\cite{hu2019local} and Ramachandran et al.~\cite{ramachandran2019stand} apply scalar dot-product self-attention within local image patches. Zhao et al.~\cite{zhao2020san} develop a family of vector self-attention operators. Dosovitskiy et al.~\cite{dosovitskiy2021image} treat images as sequences of patches.

Our work is inspired by the findings that transformers and self-attention networks can match or even outperform convolutional networks on sequences and 2D images. 
Self-attention is of particular interest in our setting because it is intrinsically a set operator: positional information is provided as attributes of elements that are processed as a set~\cite{vaswani2017attention,zhao2020san}. Since 3D point clouds are essentially sets of points with positional attributes, the self-attention mechanism seems particularly suitable to this type of data. We thus develop a Point Transformer layer that applies self-attention to 3D point clouds.

There are a number of previous works~\cite{xie2018attentional,liu2019point2sequence,yang2019modeling,lee2019set} that utilize attention for point cloud analysis. They apply global attention on the whole point cloud, which introduces heavy computation and renders these approaches inapplicable to large-scale 3D scene understanding. They also utilize scalar dot-product attention, where different channels share the same aggregation weights. In contrast, we apply self-attention locally, which enables scalability to large scenes with millions of points, and we utilize vector attention, which we show to be important for achieving high accuracy. We also demonstrate the importance of appropriate position encoding in large-scale point cloud understanding, in contrast to prior approaches that omitted position information. Overall, we show that appropriately designed self-attention networks can scale to large and complex 3D scenes, and substantially advance the state of the art in large-scale point cloud understanding.